\documentclass[12pt,a4paper]{article}

% =========================================================
% PAQUETES
% =========================================================

% Idioma
\usepackage[spanish,es-nodecimaldot]{babel}

% Codificación y fuentes
\usepackage[utf8]{inputenc}
\usepackage[T1]{fontenc}

% Matemáticas
\usepackage{amsmath}
\usepackage{amssymb}
\usepackage{amsfonts}

% Imágenes
\usepackage{graphicx}
\usepackage{float}

% Tablas
\usepackage{booktabs}
\usepackage{array}

% Márgenes
\usepackage[
    top=2.5cm,
    bottom=2.5cm,
    left=3cm,
    right=2.5cm
]{geometry}

% Interlineado
\usepackage{setspace}
\onehalfspacing

% Enlaces
\usepackage[hidelinks]{hyperref}

% Encabezados y pies de página
\usepackage{fancyhdr}

% =========================================================
% CONFIGURACIÓN
% =========================================================

\pagestyle{fancy}
\fancyhf{}

\fancyhead[L]{Informe de Ingeniería}
\fancyhead[R]{\thepage}

\renewcommand{\headrulewidth}{0.4pt}
\renewcommand{\footrulewidth}{0pt}

% =========================================================
% INFORMACIÓN DEL DOCUMENTO
% =========================================================

\title{
    \textbf{TÍTULO DEL INFORME}\\[0.5cm]
    \large Nombre de la asignatura
}

\author{
    SNEIDER YAMID SANABRIA ALFONZO\\
    Ingeniería Electrónica\\
    Nombre de la Universidad
}

\date{\today}

% =========================================================
% INICIO DEL DOCUMENTO
% =========================================================

\begin{document}

% ---------------------------------------------------------
% PORTADA
% ---------------------------------------------------------

\begin{titlepage}

    \centering

    {\Large \textbf{NOMBRE DE LA UNIVERSIDAD}\par}
    
    \vspace{0.5cm}
    
    {\large Facultad de Ingeniería\par}
    
    \vspace{0.5cm}
    
    {\large Programa de Ingeniería Electrónica\par}

    \vfill

    {\LARGE \textbf{TÍTULO DEL INFORME}\par}

    \vspace{1cm}

    {\large Informe de laboratorio / proyecto\par}

    \vfill

    \textbf{Presentado por:}\\
    SNEIDER YAMID SANABRIA ALFONZO

    \vspace{1cm}

    \textbf{Asignatura:}\\
    Nombre de la asignatura

    \vfill

    Docente: Nombre del docente

    \vspace{1cm}

    Ciudad, Colombia\\
    \today

\end{titlepage}

% ---------------------------------------------------------
% TABLA DE CONTENIDO
% ---------------------------------------------------------

\tableofcontents

\newpage

% =========================================================
% CONTENIDO
% =========================================================

\section{Introducción}

En esta sección se presenta el contexto general del proyecto,
laboratorio o trabajo realizado. Se describe brevemente el
problema abordado, la importancia del desarrollo y los
principales aspectos que serán tratados en el informe.

\section{Objetivos}

\subsection{Objetivo general}

Describir aquí el objetivo general del proyecto o práctica.

\subsection{Objetivos específicos}

\begin{itemize}
    \item Primer objetivo específico.
    \item Segundo objetivo específico.
    \item Tercer objetivo específico.
\end{itemize}

\section{Marco teórico}

En esta sección se presentan los conceptos teóricos necesarios
para comprender el desarrollo del proyecto.

\subsection{Concepto 1}

Explicación del concepto.

\subsection{Concepto 2}

Explicación del concepto.

\subsection{Ecuaciones}

Por ejemplo, una ecuación puede escribirse como:

\begin{equation}
    V = IR
    \label{eq:ohm}
\end{equation}

donde $V$ representa el voltaje, $I$ la corriente y $R$ la
resistencia.

La ecuación \ref{eq:ohm} corresponde a la Ley de Ohm.

\section{Metodología}

En esta sección se describe paso a paso el procedimiento
realizado para desarrollar el proyecto o práctica.

\subsection{Materiales y equipos}

\begin{itemize}
    \item Osciloscopio.
    \item Multímetro.
    \item Fuente de alimentación.
    \item Protoboard.
    \item Componentes electrónicos.
\end{itemize}

\subsection{Procedimiento}

Describir aquí el procedimiento realizado.

\section{Diseño e implementación}

En esta sección se presenta el diseño del sistema y su
implementación.

\subsection{Diagrama del sistema}

Para insertar una imagen:

\begin{figure}[H]
    \centering
    \includegraphics[width=0.8\textwidth]{imagenes/diagrama.png}
    \caption{Diagrama general del sistema.}
    \label{fig:diagrama}
\end{figure}

La Figura \ref{fig:diagrama} presenta el diagrama general
del sistema.

\subsection{Montaje experimental}

Describir el montaje realizado.

\section{Resultados}

En esta sección se presentan los resultados obtenidos durante
la práctica o proyecto.

\subsection{Resultados experimentales}

Los resultados obtenidos se presentan en la Tabla
\ref{tab:resultados}.

\begin{table}[H]
    \centering
    \caption{Resultados experimentales.}
    \label{tab:resultados}
    
    \begin{tabular}{ccc}
        \toprule
        Variable & Valor & Unidad \\
        \midrule
        Voltaje & 5 & V \\
        Corriente & 0.5 & A \\
        Resistencia & 10 & $\Omega$ \\
        \bottomrule
    \end{tabular}
\end{table}

\subsection{Análisis de resultados}

Analizar aquí los resultados obtenidos y compararlos con los
valores teóricos.

\section{Conclusiones}

\begin{itemize}
    \item Primera conclusión.
    \item Segunda conclusión.
    \item Tercera conclusión.
\end{itemize}

\section{Referencias}

\begin{thebibliography}{9}

\bibitem{ref1}
Autor, \textit{Título del libro o documento},
Editorial, año.

\bibitem{ref2}
Autor, ``Título del artículo'',
Nombre de la revista, año.

\end{thebibliography}

\end{document}